\chapter{Reference: CMake integration}
\label{ch:refman-cmake}

\begin{aside}
\maya{} ships a CMake package config. Find it; link to
\code{maya::maya}; you're done.
\end{aside}

\section{Minimum project}

A working project has three files: \code{CMakeLists.txt},
\code{main.cpp}, and a build directory.

\begin{lstlisting}[language=cmake]
# CMakeLists.txt
cmake_minimum_required(VERSION 3.20)
project(hello CXX)

set(CMAKE_CXX_STANDARD 23)
set(CMAKE_CXX_STANDARD_REQUIRED ON)

find_package(maya REQUIRED)

add_executable(hello main.cpp)
target_link_libraries(hello PRIVATE maya::maya)
\end{lstlisting}

\begin{lstlisting}[language=C++]
// main.cpp
#include <maya/maya.hpp>

using namespace maya;
using namespace maya::dsl;

int main() {
    print(t<"hello"> | Bold | Fg<100, 220, 160>);
}
\end{lstlisting}

\section{Build it}

\begin{lstlisting}[language=bash]
cmake -S . -B build
cmake --build build -j
./build/hello
\end{lstlisting}

\section{Linking against widgets}

Widgets live in headers; they pull no extra targets. Just
include them:

\begin{lstlisting}[language=C++]
#include <maya/widget/inline_diff.hpp>
\end{lstlisting}

\section{Vendoring instead of installing}

If you want \maya{} as a submodule:

\begin{lstlisting}[language=cmake]
add_subdirectory(third_party/maya)
target_link_libraries(hello PRIVATE maya::maya)
\end{lstlisting}

\section{Notes}

\begin{itemize}[leftmargin=*]
  \item \code{find\_package(maya REQUIRED)} fails loud if
    the package isn't installed; that's good.
  \item Use \code{target\_compile\_features(... cxx\_std\_23)}
    for older CMake versions that don't honour the global
    \code{CMAKE\_CXX\_STANDARD}.
  \item For sanitisers, add \code{-fsanitize=address,undefined}
    to \code{target\_compile\_options} in debug builds.
\end{itemize}

\section{Recap}

\begin{recap}
\begin{itemize}[leftmargin=*]
  \item \code{find\_package(maya REQUIRED)}.
  \item Link \code{maya::maya}.
  \item Widgets are header-only.
\end{itemize}
\end{recap}

\section*{Workshop}

\begin{enumerate}[leftmargin=*]
  \item Create a fresh \maya{} project from scratch in five
    minutes.
  \item Add a unit test target via \code{add\_test}.
  \item Wire up a sanitiser-only build configuration.
\end{enumerate}
